\section{自主机器人学习的逻辑脉络}
\subsection{Human-in-the-loop 与 Autonomy 的落差}
当前机器人学习体系依赖于人类的监督：奖励函数、环境重置、数据采集均靠工程师完成。《Autonomy》讲座提出挑战：让机器人在复杂物理世界中像婴儿一样主动探索，而不是等着人类定目标。

\begin{figure}[H]
\centering
\includegraphics[width=0.92\textwidth]{slides-images/slide-01.jpg}
\caption{Slide 1 概览讲座的核心主题：Autonomy 既是 reward, reset, goal 设置的集合。}
\end{figure}

\begin{importantbox}{Autonomy 的三维目标}
1. 自动构造奖励或理解任务成果；\\
2. 在不重置的世界里持续操作；\\
3. 自主选择目标并安全探索。
\end{importantbox}

\subsection{面对现实的评估标准}
讲者强调三个评价维度：\textbf{可学习性}（是否能获得足够 signal）、\textbf{可扩张性}（是否不依赖人工重置）、\textbf{可持续性}（是否能长期运行并在新的目标上迁移）。教学逻辑：先从 reward 说起，再从环境与目标的逻辑演化到安全与系统集成。

\begin{knowledgebox}{教学逻辑映射}
Autonomy 研究线索依次是：\textbf{自动奖励}→\textbf{Reset-free 控制}→\textbf{Goal Management}→\textbf{Safe Exploration}→\textbf{系统监控}。每一段都以旧假设（reward, reset, goal）为起点，描述如何撤回它们。
\end{knowledgebox}

\subsection{本章小结}
本章梳理了自主机器人学习的现状与教学逻辑，以“奖励-环境-目标-安全-系统”五阶段为后续章节的纲领框架。
\newpage

